\FloatBarrier

\section{Electron and hadronic transport}

Table~\ref{tab:3} gives median events-only CPU time per 1,000 primaries. The nonoptical entries divide the
measured 10,000-event times by ten. Throughput factors are Original time divided by comparison time: values
above one indicate faster transport, and values below one indicate slower transport. Figure~\ref{fig:transport}
presents these ratios and the initialization ratio for direct comparison.

\begin{table}[htbp]

\caption{CPU seconds per 1,000 events and transport throughput factors; zero magnetic field.}

\label{tab:3}

\centering

\small

\setlength{\tabcolsep}{4pt}

\renewcommand{\arraystretch}{1.15}

\begin{tabularx}{\textwidth}{@{}Xrrrrr@{}}

\toprule

Workload & Original & Cones only & Full port & Cone factor & Port factor \\

\midrule

Electron + optical & 26.067 & 26.852 & 45.126 & 0.97 & 0.58 \\

Electron & 4.739 & 4.546 & 3.679 & 1.04 & 1.29 \\

Proton & 1.009 & 0.958 & 1.490 & 1.05 & 0.68 \\

Positive pion & 1.084 & 1.055 & 1.467 & 1.03 & 0.74 \\

Neutron & 0.575 & 0.601 & 0.877 & 0.96 & 0.66 \\

\bottomrule

\end{tabularx}

\end{table}

Optical electron transport: the cone-only substitution has a throughput factor of 0.97; the full port has a
factor of 0.58. Original per-1,000-event CPU times range from 25.549 to 29.228 s, and full-port times range
from 39.580 to 56.916 s.

Nonoptical electron transport: the cone-only substitution has a throughput factor of 1.04; the full port has a
factor of 1.29. Original per-1,000-event CPU times range from 4.207 to 6.261 s, and full-port times range from
3.476 to 6.536 s.

Proton transport: the cone-only substitution has a throughput factor of 1.05; the full port has a factor of
0.68. Original per-1,000-event CPU times range from 0.987 to 1.034 s, and full-port times range from 1.477 to
1.580 s.

Positive-pion transport: the cone-only substitution has a throughput factor of 1.03; the full port has a factor
of 0.74. Original per-1,000-event CPU times range from 1.060 to 1.170 s, and full-port times range from 1.444
to 1.546 s.

Neutron transport: the cone-only substitution has a throughput factor of 0.96; the full port has a factor of
0.66. Original per-1,000-event CPU times range from 0.543 to 0.651 s, and full-port times range from 0.842 to
0.908 s.

The cone-only comparison isolates the three substituted meshes. The full port additionally changes gas
containment, hardware coverage, placements, and materials. Optical paths and particle showers can therefore
differ, and faster loading does not imply faster transport or unchanged response. The measured changes include
both mesh-navigation cost and changes in physical transport; these data do not identify a unique cause. PMT
absorption was held fixed and cannot account for the difference between these cases.

Cone-only throughput ratios close to one represent small measured differences within the observed run
variation; these five repetitions do not establish a statistically significant speedup of the Winston-cone
substitution. The full-model timing also includes the revised materials and mothers, and does not separately
measure CPU cost per navigation query.

A separate matched 100-event electron diagnostic at seed 12345 enables tracking verbosity 1 and counts the
printed trace (Original: 18,271 track headers and 123,928 step rows; Cones only: 17,025 track headers and
115,895 step rows; Full port: 24,164 track headers and 203,087 step rows). Step rows include the initial row
for each track; they are not a direct count of navigation queries. These counts characterize the amount of
transport work in this sample. Verbose-run times are excluded from the timing tables, and this sample does not
isolate the contribution of each geometry change.

A separate matched 100-event proton diagnostic at seed 12345 enables tracking verbosity 1 and counts the
printed trace (Original: 6,135 track headers and 50,512 step rows; Cones only: 7,794 track headers and 62,163
step rows; Full port: 9,566 track headers and 99,568 step rows). Step rows include the initial row for each
track; they are not a direct count of navigation queries. These counts characterize the amount of transport
work in this sample. Verbose-run times are excluded from the timing tables, and this sample does not isolate
the contribution of each geometry change.

The increased CPU time is consistent with more complete transport in the optimized geometry. The original model
has invalid or rejected facets and repeated stuck tracks requiring corrective navigation pushes; these defects
may alter material crossings and suppress expected particle histories. In the matched diagnostic samples, the
optimized model prints 1.64 times as many electron step rows and 1.97 times as many proton step rows. Lower CPU
time in the defective original model may therefore reflect incorrect or incomplete transport rather than a
lower cost of equivalent physics. These observations support that interpretation but do not establish
causation: materials and mothers also change, and this study does not compare energy balance, particle
survival, or detector response event by event.

\begin{figure}[htbp]

\centering

\includegraphics[width=\textwidth]{figures/transport_comparison.pdf}

\caption{Original/comparison CPU-time ratios for initialization and particle transport. Bars use the ratio of
five-run medians. Transport uses the isolated default LTCC at RUNNO=11, 5 GeV/c primary momentum, and zero
magnetic field; only the optical electron workload enables optical processes. No readout threshold or
luminosity normalization applies. Ratios above one mean faster.}

\label{fig:transport}

\end{figure}
